\subsection{The complete level-three models for e = 20, 22, and 24}
\label{sec:u3-higher3}

\begin{proposition}[K17-HECKE-HIGHER3]
\label{prop:hecke-higher3}
Retain the full compact-unitary and classical models of
Theorem~\ref{thm:u3-general} and Proposition~\ref{prop:gl2-general}.
For \(e=20,22,24\) and \(p=7,13\), let \(P_{e,p}\) be the
characteristic polynomial of \(p^{e-1}T_p^+\) on the entire
compact-unitary coefficient space of highest weight
\((e-1,0,1-e)\), and let \(Q_{e,p}\) be the characteristic
polynomial of the usual \(T_p\) on
\(S_{2e+1}(\Gamma_0(3),\chi_{-3})\).
The complete dimensions are
\[
\begin{array}{c|c|r|r|r|c|r}
e&k=e-1&\dim V_k^\Gamma&\dim V_{k-1}^\Gamma&
\dim\mathcal H_k^\Gamma&\text{swap }(+,-)&
\dim S_{2e+1}\\ \hline
20&19&230&191&39&(21,18)&12\\
22&21&329&277&52&(28,24)&14\\
24&23&452&388&64&(34,30)&15 .
\end{array}
\tag{HH3.1}
\]
Put \(E_{e,p}=1+p^e+p^{2e}\) and
\[
 N_{e,p}=P_{e,p}(E_{e,p})Q_{e,p}(1+p^{2e}).
 \tag{HH3.2}
\]
Define the following positive integers:
\[
\begin{array}{c|r}
M_{20,1}&9496386083\\
M_{20,2}&552972895573825165951583191\\
M_{22,1}&80348736972167\\
M_{22,2}&2789260013036301774436972122687974510453\\
M_{24,1}&3369865772726176603506574274451733385689247 .
\end{array}
\tag{HH3.3}
\]
The exact common divisors are
\[
\begin{aligned}
\gcd(N_{20,7},N_{20,13})
 &=2^{213}3^{257}5^6\,61^2\,241\,175427911^3\,
        M_{20,1}M_{20,2}^3,\\
\gcd(N_{22,7},N_{22,13})
 &=2^{298}3^{299}5^{29}\,19^2\,53^3\,67\,
        M_{22,1}M_{22,2}^3,\\
\gcd(N_{24,7},N_{24,13})
 &=2^{298}3^{372}5^8\,7\,13^9\,37\,97\,28447\,
        300455201\,M_{24,1}^3 .
\end{aligned}
\tag{HH3.4}
\]
Every displayed factor other than an \(M_{e,i}\) is a proved
prime below \(2^{31}\).
The numbers \(M_{e,i}\) are residual integers; no assertion
of primality is made or needed.
\end{proposition}

\begin{proof}
The full-space argument in Theorem~\ref{thm:u3-general} is
uniform in \(k\). In bidegree \((k,k)\), invariance under the
diagonal sixth roots of unity is exactly
\[
 |\alpha|=|\beta|=k,\qquad
 \alpha_i-\beta_i\equiv0\pmod6\quad(i=1,2,3).
 \tag{HH3.5}
\]
Taking every simultaneous permutation orbit gives the entire
invariant space. In the present degrees this includes exponent
differences \(0,\pm6,\pm12,\pm18\); no such orbit is discarded.
The exact enumeration gives the source and target counts in
(HH3.1). The independent replay reconstructs the full contraction
matrix by differentiation of every orbit sum, proves its full
target rank, verifies that every saved harmonic column is in
its kernel, and checks independence of those columns. The
dimension difference proves completeness of the saved harmonic
basis. Swapping the two variable triples gives the recorded
plus and minus dimensions.

The same complete sets of \(57\) and \(183\) neighbors at \(7\)
and \(13\) are used. The replay verifies their exact Hermitian
Gram matrices, their distinct residue lines, and the equality
of those lines with every point of \(\mathbf P^2(\mathbf F_p)\).
For every basis vector \(h\) and every neighbor \(U/\pi\),
the scaled action is the uniform identity
\[
 p^{e-1}\rho(U/\pi)h(z,w)=h(U^*z,U^tw).
 \tag{HH3.6}
\]
The discovery code selects evaluation points by modular rank
screening and then proves the exact evaluation determinant
nonzero over \(\mathbf Q\). The independent replay reconstructs
this entire integer evaluation matrix and its determinant.
It recomputes every transformed evaluation at every neighbor
and checks both coefficient equations \(AM_0=B_0\) and
\(AM_1=B_1\). Therefore the complete two-component matrix over
\(\mathbf Q(\sqrt{-3})\) is the matrix of the full operator.
It also checks every column at two further point pairs.

The variable swap exchanges the two minuscule directions.
Multiplying the swap-minus basis by \(1+2\omega\) produces
the same rational form as in the earlier model theorem; the
replay checks that change of form entry by entry. The scaled
operator preserves the full lattice of invariant harmonic
polynomials over \(\mathbf Z[\omega]\), so its rational
characteristic polynomial has integer coefficients.
Independently, the replay calculates the whole polynomial
using exact fractional Faddeev--LeVerrier recursion and
checks its terminal Cayley--Hamilton zero matrix. It verifies
the saved exact product decompositions, whose degree patterns
at both auxiliary primes are \(3+6+30\), \(3+7+42\), and
\(3+7+54\), respectively. Only multiplication of these factors
is used; irreducibility is not assumed.

The full classical space has the integral Fourier basis
\[
 Da^{\,2e-5-3j}b^j,\qquad
 0\le j\le\left\lfloor(2e-5)/3\right\rfloor ,
 \tag{HH3.7}
\]
by Proposition~\ref{prop:gl2-general}. Here \(D,a,b\) have
the fixed level-three meanings used there. Their successive
leading orders give dimensions \(12,14,15\) and a unit
triangular identifying Fourier matrix. The new classical
certificate reconstructs every series to precision \(260\),
which exceeds \(13\cdot15=195\), and constructs both complete
Hecke matrices from the exact formula
\[
 a_n(T_pf)=a_{pn}(f)+
 \begin{cases}
   \chi_{-3}(p)p^{2e}a_{n/p}(f),&p\mid n,\\
   0,&p\nmid n.
 \end{cases}
 \tag{HH3.8}
\]
All corrections, including those at identifying indices
divisible by \(7\) or \(13\), are retained.
It checks every coefficient through \(37\) at \(7\) and
\(20\) at \(13\). Exact Faddeev--LeVerrier recursion and
three independent fraction-free determinant evaluations,
at \(0,1,1+p^{2e}\), verify every classical polynomial.

The complete unitary data, classical matrices, and replay
reports are in
\path{hecke_general/higher3/harmonic_hecke_exact_data.json},
\path{hecke_general/higher3/gl2_higher3_certificate.json},
and the corresponding certificate reports.
The unchanged independent unitary verifier is
\path{hecke_general/harmonic_hecke_certificate.py}, run
with explicit paths to this new data and report.
It makes no primality assertions about the residual integers.

Finally,
\path{hecke_general/higher3/combined_higher3_certificate.py}
evaluates all twelve complete characteristic polynomials
by integer Horner arithmetic, checks every saved evaluation,
and computes the two-value gcds in (HH3.4) by the Euclidean
algorithm. It proves each claimed small prime by trial division
through its integer square root. Exact multiplication of those
small prime powers and the stated residual integer powers
equals the entire gcd in every case, with no unreported
cofactor. The largest trial-division bound is
\(\lfloor\sqrt{300455201}\rfloor=17333\).
The residual entries are checked only as integers and are
left intact for the separate level-seven unit tests.
\end{proof}

\paragraph{Scope of the comparison.}
Under the separately established admissibility, conductor,
global transfer, and trace interface, a coefficient prime
greater than \(2^{31}\) which does not divide any corresponding
\(M_{e,i}\) is eliminated by at least one of \(p=7,13\).
This follows from (HH3.4) and the direct unit criterion.
The proposition alone makes no arithmetic assertion at
a prime divisor of one of the residual integers.
